\documentclass[11pt]{article}
\usepackage[margin=0.8in]{geometry}
\usepackage[T1]{fontenc}
\usepackage{lmodern}
\usepackage{xcolor}
\usepackage{hyperref}
\usepackage{graphicx}
\usepackage{booktabs}
\usepackage{longtable}
\usepackage{array}
\usepackage{listings}
\usepackage{caption}
\usepackage{amsmath,amssymb,bm}

\hypersetup{
  colorlinks=true,
  linkcolor=blue!50!black,
  urlcolor=blue!50!black,
  citecolor=blue!50!black
}

\lstset{
  basicstyle=\ttfamily\scriptsize,
  breaklines=true,
  breakatwhitespace=false,
  columns=fullflexible,
  frame=single,
  rulecolor=\color{black!20},
  xleftmargin=0.5em,
  xrightmargin=0.5em
}

\newcommand{\artifact}[1]{\path{#1}}

\title{\textbf{Complete Master Compendium}\\
Controller-Aware Spectral Graph Theory for Damping Margins in IBR Grids}
\author{Internal research consolidation}
\date{Compiled \today}

\begin{document}
\maketitle

\begin{abstract}
This compendium is intentionally redundant. It preserves the full IEEE
manuscript, the unified theory note with the Gemini extensions, the diagrams,
and the validation reports. This front matter explains how the final merged PDF
is organized and appends the raw gate ledgers. The source papers themselves are
merged after this front matter as complete PDFs, so their equations and figures
are not rewritten or dropped.
\end{abstract}

\tableofcontents
\clearpage

\section*{How to Read This Compendium}
\addcontentsline{toc}{section}{How to Read This Compendium}

The final merged artifact has four layers:
\begin{enumerate}
\item this front matter and artifact appendix;
\item the unified theory note with calibrated Gemini additions;
\item the full IEEE Transactions manuscript with all equations, tables, figures,
derivation appendices, gates, and reviewer checklists;
\item raw validation reports and compact CSV ledgers.
\end{enumerate}

The document is deliberately not concise. Its role is to preserve the entire
research state before splitting the material into submission-sized papers.

\section*{Recovered Content Manifest}
\addcontentsline{toc}{section}{Recovered Content Manifest}

\begin{longtable}{p{0.30\linewidth}p{0.62\linewidth}}
\toprule
Source & Preserved content\\
\midrule
\artifact{internal\_unified\_theory/main.pdf}
& Gemini extensions after correction: protected resolvent sensitivity,
flow-feasible retrofitting candidate, damping-operator hierarchy, and explicit
novelty boundaries.\\
\artifact{paper\_ieee\_transactions/main.pdf}
& Full prior manuscript: all equations, the second-order correction, rational
graph-filter identity, joint inertia--damping perturbation, ANDES IEEE39
sections, NEP bridge, Braess mechanism, hidden margin, weak nodes/links,
planning scores, line reinforcement, appendices, and all manuscript figures.\\
\artifact{outputs/phase0d\_nep\_contour\_solver}
& Non-circular Beyn/NEP bridge validation against ANDES, including Mix60
control-family critical mode at 13.742 Hz.\\
\artifact{outputs/phase\_braess\_nep}
& Resonant Braess mechanism test and classification: strict resonant minority
cases plus negative controls.\\
\artifact{outputs/phase\_inertia\_bridge\_gate}
& Gate for inertia sensitivity over the controller-aware bridge.\\
\artifact{outputs/phase\_coordination\_lever\_gate}
& Reduced-model gate for coordinated topology/damping/inertia packages and the
absence of robust nonlinear synergy.\\
\artifact{outputs/phase\_gamma\_robustness\_gate}
& Robustness of the protected-resolvent discrepancy under physical and random
diagonal protection profiles.\\
\artifact{outputs/final\_claim\_audit}
& Final claim-readiness ledger and status labels.\\
\bottomrule
\end{longtable}

\clearpage
\section{Explicit Gemini Addendum: Flow-Feasible Graph-Filter Shaping}

The newest Gemini proposal is worth preserving because it points to a clean
conference-paper split:
\begin{quote}
\emph{Flow-Feasible Graph-Filter Shaping: Coordinating Topology, Damping, and
Virtual Inertia in Non-Normal IBR Grids.}
\end{quote}
The proposal is included here, but with the same claim discipline used
throughout the project.  The useful idea is not that every formula in the raw
proposal is already a theorem.  The useful idea is that the local protected
resolvent provides a common currency for topology, damping, and virtual inertia
when the state matrix is non-normal.

\subsection{Gemini G1: Bi-Orthogonal Resolvent Sensitivity}

Let
\[
    H(j\omega)=C(j\omega I-A_{\rm sys})^{-1}B,
    \qquad
    g_{\rm prot}=\sigma_1(H(j\omega^\star)),
\]
where $\omega^\star$ is the protected-resonance peak.  Let $u^\star$ and
$v^\star$ be the left and right singular vectors of $H(j\omega^\star)$ associated
with $g_{\rm prot}$.  Define
\[
    \phi^{\rm in}=(j\omega^\star I-A_{\rm sys})^{-1}Bv^\star ,
    \qquad
    \phi^{\rm out}=(j\omega^\star I-A_{\rm sys})^{-H}C^H u^\star .
\]
For a smooth real parameter $p$ in $A_{\rm sys}$, the first-order sensitivity of
the protected log-gain $\mathcal S=\log g_{\rm prot}$ is
\[
    \boxed{
    \frac{\partial \mathcal S}{\partial p}
    =
    \frac{1}{g_{\rm prot}}
    \operatorname{Re}\!\left[
      (\phi^{\rm out})^H
      \frac{\partial A_{\rm sys}}{\partial p}
      \phi^{\rm in}
    \right].}
\]

\textbf{Status.} This is standard singular-value perturbation applied to the
power-system resolvent.  The contribution is the use of this expression as a
planning exchange rate across grid levers in an IBR damping-margin setting, not
the invention of singular-value perturbation.

\subsection{Gemini G2: Spatial Decoupling, Corrected}

The raw Gemini note states a strong damping-placement rule of the form
$\partial \mathcal S/\partial D_i\propto-|\phi^{\rm out}_i|^2$.  That rule is
\emph{not} generally valid for non-normal systems.  In the general case the
correct local damping sensitivity is a bi-orthogonal product.  For a second-order
state $x=[\theta^T,\omega^T]^T$ with local damping entering the $\omega_i$
equation as $-\Delta D_i/M_i$, one obtains
\[
    \boxed{
    \frac{\partial \mathcal S}{\partial D_i}
    =
    -\frac{1}{g_{\rm prot}M_i}
    \operatorname{Re}\!\left[
      \overline{\phi^{\rm out}_{\omega_i}}
      \phi^{\rm in}_{\omega_i}
    \right].}
\]
The simplified $-|\phi_i|^2$ expression is recovered only in aligned or nearly
normal cases where input and output spatial profiles coincide.  The actual
non-normal insight is stronger:
\begin{quote}
the location that receives forcing and the location where response should be
damped need not be the same bus.
\end{quote}

For a line $e=(i,j)$ with incidence vector $b_e=e_i-e_j$ and stiffness change
$w_e$, the corresponding topology sensitivity has the form
\[
    \boxed{
    \frac{\partial \mathcal S}{\partial w_e}
    =
    -\frac{1}{g_{\rm prot}}
    \operatorname{Re}\!\left[
      (\phi^{\rm out}_{\omega})^H
      M^{-1} b_e b_e^T
      \phi^{\rm in}_{\theta}
    \right],}
\]
up to the precise state ordering and sign convention used in the assembled
$A_{\rm sys}$.  This is the graph-filter version of a weak-link score: a line is
dangerous or useful according to how it couples the forcing profile to the
response profile, not only according to Fiedler participation or stiffness gain.

\textbf{Status.} The certificate-discrepancy gates support that the modal damping
margin $S_\zeta$ and protected resolvent $S_g$ disagree robustly under several
protection profiles.  The mechanism of the discrepancy is not fully explained:
global non-normality and global commutator descriptors were weak predictors.
The spatial-decoupling theorem is therefore a strong candidate, not yet a closed
validated theorem over ANDES.

\subsection{Gemini G3: Flow-Feasible Dynamic Braess}

A line reinforcement is a \emph{flow-feasible dynamic Braess} candidate if, at a
given operating point, it satisfies the three local conditions
\[
    \frac{\partial \nu_c}{\partial w_e}>0,
    \qquad
    \frac{\partial \mathcal S}{\partial w_e}>0,
    \qquad
    |f_0+\Delta f(w_e,P)|\leq \bar f .
\]
The first condition says that the line improves a frequency or stiffness proxy.
The second says that the protected dynamic gain worsens.  The third prevents a
purely dynamic claim from ignoring steady-state flow feasibility.

The local retrofit problem suggested by Gemini can be written as
\[
    \begin{aligned}
    \min_{z=(\Delta w,\Delta D,\Delta M)}\quad & c^Tz\\
    \text{s.t.}\quad
       & \mathcal S_0+\nabla\mathcal S^Tz \leq \mathcal S_{\rm target},\\
       & F_0+J_F z \leq \bar F,\\
       & z_{\min}\leq z\leq z_{\max}.
    \end{aligned}
\]
This is not yet a global optimality theorem.  It is a local planning layer that
can be tested against full NEP or ANDES eigensolves after a candidate package is
selected.

\subsection{Algorithm G1: Flow-Feasible Resolvent Retrofitting}

\begin{enumerate}
\item Build or import $A_{\rm sys}$, protected input/output matrices $B,C$, flow
limits $\bar F$, and action costs $c$.
\item Find $\omega^\star$, $g_{\rm prot}$, and singular vectors $u^\star,v^\star$
of $H(j\omega)$.
\item Compute $\phi^{\rm in}$ and $\phi^{\rm out}$ from the two resolvent solves.
\item Generate exchange rates
$\partial\mathcal S/\partial w_e$,
$\partial\mathcal S/\partial D_i$, and
$\partial\mathcal S/\partial M_i$.
\item Compute flow penalties from $J_F$ or from an AC/DC sensitivity model.
\item Solve the local retrofit LP for a sparse package
$z^\star=(\Delta w^\star,\Delta D^\star,\Delta M^\star)$.
\item Validate the selected package against the NEP/ANDES poles.  If the
critical singular vectors or pole family switch, rerun the local model.
\end{enumerate}

\subsection{Figures Needed for This Spin-Off Paper}

The Gemini note proposes three figures.  They should be generated only after the
corresponding gates are passed.
\begin{enumerate}
\item \textbf{Spatial decoupling heatmap:} forcing bus/location versus damping
placement bus/location from $\phi^{\rm in}$ and $\phi^{\rm out}$.
\item \textbf{Flow-feasible Braess witness:} a line that improves the stiffness
proxy, keeps flows legal, but worsens the protected resolvent or damping margin.
\item \textbf{Cost-synergy frontier:} topology-only, damping/inertia-only, and
coordinated packages.  Existing gates support $w+D$ and $w+M$ as reduced-model
coordination candidates, but do not support a strong nonlinear synergy claim.
\end{enumerate}

\subsection{Claim Boundary for the Gemini Direction}

This direction is promising because it converts weak nodes and weak links from
single-score graph objects into input-output objects.  The claim should be:
\begin{quote}
Protected-resolvent sensitivities define a common local exchange rate for
topology, damping, and inertia actions in non-normal IBR grids, and early
reduced-model gates show robust disagreement with modal damping-margin rankings.
\end{quote}
It should \emph{not} yet be:
\begin{quote}
The optimal damping bus is always the output hot spot, or coordinated packages
are always synergistic, or the LP is globally optimal.
\end{quote}

\clearpage
\section{Figure and Diagram Recovery Appendix}

\subsection{Manuscript Figures}

\begin{figure}[h]
\centering
\includegraphics[width=0.82\linewidth]{paper_ieee_transactions/figures/fig1_damping_region.pdf}
\caption{Damping-region geometry from the full manuscript.}
\end{figure}

\begin{figure}[h]
\centering
\includegraphics[width=0.82\linewidth]{paper_ieee_transactions/figures/fig2_estimator_errors.pdf}
\caption{Estimator error stress test from the full manuscript.}
\end{figure}

\begin{figure}[h]
\centering
\includegraphics[width=0.82\linewidth]{paper_ieee_transactions/figures/fig3_adaptive_workflow.pdf}
\caption{Adaptive workflow: diagonal screen, correction, and fallback.}
\end{figure}

\begin{figure}[h]
\centering
\includegraphics[width=0.82\linewidth]{paper_ieee_transactions/figures/fig4_line_tradeoff.pdf}
\caption{Frequency-only versus damping-aware line reinforcement tradeoff.}
\end{figure}

\begin{figure}[h]
\centering
\includegraphics[width=0.92\linewidth]{paper_ieee_transactions/figures/fig5_ieee39_ibr_map.pdf}
\caption{IEEE39 SG-to-IBR replacement map.}
\end{figure}

\begin{figure}[h]
\centering
\includegraphics[width=0.92\linewidth]{paper_ieee_transactions/figures/fig6_ieee39_result_map.pdf}
\caption{IEEE39 result map and validation status overview.}
\end{figure}

\begin{figure}[h]
\centering
\includegraphics[width=0.82\linewidth]{paper_ieee_transactions/figures/fig7_braess_nep_s_plane.pdf}
\caption{Control-resonant Braess example in the complex plane.}
\end{figure}

\clearpage
\subsection{Output Figures}

\begin{figure}[h]
\centering
\includegraphics[width=0.92\linewidth]{outputs/phase0d_nep_contour_solver/fig_phase0d_s_plane.pdf}
\caption{Phase-0D NEP contour validation: ANDES poles, NEP zeros, control poles,
and contours.}
\end{figure}

\begin{figure}[h]
\centering
\includegraphics[width=0.82\linewidth]{outputs/phase_braess_nep/fig_braess_nep_s_plane.pdf}
\caption{Phase Braess-NEP resonant line-reinforcement witness.}
\end{figure}

\begin{figure}[h]
\centering
\includegraphics[width=0.92\linewidth]{outputs/ieee39_phase0_bridge/fig_ieee39_mix60_ibr_map.pdf}
\caption{Mix60 IBR map generated during the ANDES bridge phase.}
\end{figure}

\begin{figure}[h]
\centering
\includegraphics[width=0.92\linewidth]{outputs/ieee39_result_map/fig_ieee39_result_map.pdf}
\caption{IEEE39 result-map artifact.}
\end{figure}

\begin{figure}[h]
\centering
\includegraphics[width=0.82\linewidth]{outputs/ieee39_rational_filter/fig_estimator_errors_by_rho.pdf}
\caption{Estimator errors by IBR penetration in the rational-filter artifact set.}
\end{figure}

\begin{figure}[h]
\centering
\includegraphics[width=0.82\linewidth]{outputs/ieee39_rational_filter/fig_line_frequency_vs_damping.pdf}
\caption{Line frequency score versus damping score in the rational-filter artifact set.}
\end{figure}

\clearpage
\section{Raw Gate Reports and Result Ledgers}

\subsection{Phase-0D NEP Contour Solver Report}
\lstinputlisting{outputs/phase0d_nep_contour_solver/phase0d_nep_report.md}

\subsection{Phase Braess-NEP Report}
\lstinputlisting{outputs/phase_braess_nep/phase_braess_nep_report.md}

\subsection{Phase Inertia Bridge Gate Report}
\lstinputlisting{outputs/phase_inertia_bridge_gate/phase_inertia_bridge_gate_report.md}

\subsection{Phase Coordination Lever Gate Report}
\lstinputlisting{outputs/phase_coordination_lever_gate/phase_coordination_lever_gate_report.md}

\subsection{Phase Gamma Robustness Gate Report}
\lstinputlisting{outputs/phase_gamma_robustness_gate/phase_gamma_robustness_gate_report.md}

\subsection{Final Claim Readiness Report}
\lstinputlisting{outputs/final_claim_audit/final_readiness_report.md}

\clearpage
\newgeometry{margin=0.35in}
\section{Compact CSV Ledgers}

\lstset{basicstyle=\ttfamily\tiny}

\subsection{Claim Audit CSV}
\lstinputlisting{outputs/final_claim_audit/claim_audit.csv}

\subsection{Phase-0D Critical Validation CSV: No-PSS}
\lstinputlisting{outputs/phase0d_nep_contour_solver/no_pss_critical_validation.csv}

\subsection{Phase-0D Critical Validation CSV: Base}
\lstinputlisting{outputs/phase0d_nep_contour_solver/base_critical_validation.csv}

\subsection{Phase-0D Critical Validation CSV: Mix60}
\lstinputlisting{outputs/phase0d_nep_contour_solver/mix60_critical_validation.csv}

\subsection{Coordination Aggregate Table}
\lstinputlisting{outputs/phase_coordination_lever_gate/coordination_aggregate_table.csv}

\subsection{Gamma Aggregate Table}
\lstinputlisting{outputs/phase_gamma_robustness_gate/gamma_aggregate_table.csv}

\end{document}
